\cotefile{79}
\scanpage{265}
\mspage{135}

En résumé :

\[
\left\{\begin{aligned}
& T_0\ \text{opère fidèlement sur }H,\ \text{et via cette}\\
&\text{opération s'identifie au sous-groupe }\inserted{\widetilde Z}\ \text{de }\operatorname{Aut}(H)\\
&\text{formé des }\bar g\in\operatorname{Aut} H\ \text{satisfaisant les conditions}\\
&\qquad\left\{\begin{aligned}
&\bar g\ \text{normalise }U_0\ \text{et centralise }T_H\\
&\bar g(\sigma)\in\sigma T_H.
\end{aligned}\right.\\
&\text{De plus, on aura alors}\\
&\qquad\bar g(\sigma)=\sigma\,u(\underset{\substack{\shortparallel\\c(\widetilde\delta(g))}}{\underbrace{\delta(g)}}).
\end{aligned}\right.
\tag*{(144)}
\]

% slanted
\margin{\uline{NB} Ces conditions impliquent aussi que $\bar g$ normalise $U_1$ -- en fait, si $\bar g$ \deleted{applique} applique $\sigma T_H=T_H\sigma$ dans lui-\unsure{m\textsuperscript{m}}, alors il normalise $U_0$ ssi il normalise $U_1$.}

En d'autres termes, \deleted{si on désigne par $\widetilde Z$ une ``copie conforme'' de $T_0\subset\operatorname{Aut}(H)$, via $u_0:\widetilde Z\xrightarrow{\sim}T_0$ et $\widetilde\delta_{\!T_0}\deleted{\unsure{u_0}}\unsure{p_0}:T_0\xrightarrow{\sim}\widetilde Z$} alors l'hom $c:\widetilde Z\longrightarrow Z$ s'identifie à l'hom

\[
\deleted{T_0}\widetilde Z\xrightarrow{c_0}T_H\tag*{(145)}
\]

défini par

\[
\bar g\unsure{{}_0|{}}(\sigma)=\sigma\,c_0(\bar g).
\]

i.e on a commutativité dans

\diagram{\makebox[0pt][r]{$(145')$\quad}\begin{tikzcd}[ampersand replacement=\&,column sep=large]
\widetilde Z \arrow[r,"c"] \arrow[dr,"c_0"] \arrow[d,"\deleted{u_0}"'] \& Z \arrow[d,"u","{\unsure{?|\sim|\wr}}"']\\
\deleted{T_0} \arrow[r,"\deleted{\unsure{c_0}}"'] \& T_H
\end{tikzcd}}

Notons d'autre part que l'hom

\[
\begin{aligned}
T_H&\longrightarrow\operatorname{Aut}(H)\\
t&\longmapsto\operatorname{int}_H(t)
\end{aligned}
\]

se factorise visiblement par $\deleted{T_0}\widetilde Z$, et donne un homomorphisme

\[
T_H\xrightarrow{c'_0}\deleted{T_0}\widetilde Z\tag*{(146)}
\]

satisfaisant

\[
c_0(c'_0(t))=t^{-2}\tag*{\deleted{\unsure{(}}(147)}
\]

puisque pour $t\in T_H$, on a

\[
\sigma t\underset{\sigma}{\underbrace{\sigma^{-1}}}=t^{-1}\quad\deleted{\unsure{\text{d'où}}}\ \text{i.e.}\quad t\sigma t^{-1}=\sigma t^{-2},
\]

ce qui suggère qu'on a commutativité dans
% archive sheet number: 264
